\documentclass[11pt]{amsart}
\usepackage{amsmath,amssymb,amsthm,booktabs}
\usepackage[margin=1.15in]{geometry}
\usepackage{hyperref}
\newtheorem{theorem}{Theorem}[section]
\newtheorem{proposition}[theorem]{Proposition}
\newtheorem{lemma}[theorem]{Lemma}
\newtheorem{corollary}[theorem]{Corollary}
\theoremstyle{definition}
\newtheorem{definition}[theorem]{Definition}
\newtheorem{remark}[theorem]{Remark}
\newtheorem{example}[theorem]{Example}
\newcommand{\Prob}{\mathbb{P}}
\newcommand{\E}{\mathbb{E}}
\newcommand{\R}{\mathbb{R}}
\newcommand{\Z}{\mathbb{Z}}
\newcommand{\cx}{\le_{\mathrm{cx}}}
\newcommand{\icx}{\le_{\mathrm{icx}}}
\newcommand{\eqd}{\stackrel{d}{=}}
\newcommand{\one}{\mathbf{1}}
\newcommand{\cN}{\mathcal{N}}

\title[Norm inequalities and complete mixability]{Law-determined norm inequalities do not imply complete mixability}
\author{Yitzchak Shmalo}
\address{Einstein Institute of Mathematics, The Hebrew University of Jerusalem, Jerusalem, Israel}
\date{November 2026}
\subjclass[2020]{60E15, 60E05, 91G70}
\keywords{complete mixability, convex order, law-determined norm, rearrangement-invariant norm, majorization}

\begin{document}
\begin{abstract}
A law $F$ with mean $\mu$ is $n$-completely mixable if there are $n$ random variables with law $F$ whose sum is
almost surely $n\mu$. Wang proved that such a law satisfies a family of inequalities between the norms of the
positive and negative parts of $X-t$ and $X-s$, where $X\sim F$ and $s=(n\mu-t)/(n-1)$, for every threshold $t$ and
every law-determined norm, and asked when these inequalities are also sufficient (Open Problem~4 in his 2015
survey). We show that for bounded $F$ the whole family is equivalent to the single convex-order condition
$X-\mu\cx-(n-1)(X-\mu)$, and that the same holds for integrable $F$ if the norms are required to have the Fatou
property. The inequalities do not imply $n$-complete mixability: for $n=3$ the analytic density
$\frac1{10}+\frac{27}{10}(2x-1)^2$ on $[0,1]$ satisfies all of them and is not $3$-completely mixable, and for $n=3$
and $n=4$ there are counterexamples with densities bounded away from zero on $[0,1]$ and with $C^\infty$ densities,
the one for $n=4$ not symmetric. A Bernoulli$(p)$ law satisfies the inequalities exactly when $1/n\le p\le1-1/n$,
while it is $n$-completely mixable exactly when $np$ is an integer. For $n=2$ the inequalities are equivalent to
symmetry, hence to $2$-complete mixability, for every integrable law. For $n=3$ the examples refine the symmetric
non-mixable laws of Yin and Zhu, and the reduction explains why those laws satisfy the inequalities.
\end{abstract}
\maketitle

\section{Introduction}\label{sec:intro}

A law $F$ on $\R$ with finite mean $\mu$ is \emph{$n$-completely mixable} ($n$-CM) if there are random variables
$X_1,\dots,X_n$, each with law $F$, such that $X_1+\dots+X_n=n\mu$ almost surely. Complete mixability appears in
bounds on the value-at-risk of a sum with given marginals and in variance minimization; the survey of
Wang~\cite{W15} collects what is known and lists open problems. All random variables below live on an atomless
probability space, and $L^0$ denotes the space of all of them, modulo almost sure equality.

Complete mixability was developed as a property of laws by Bin Wang and Ruodu Wang~\cite{WW11}, after earlier
constant-sum constructions such as the symmetric unimodal one of R\"uschendorf and Uckelmann~\cite{RU02}. Ruodu Wang's
norm inequalities and his Open Problem~4~\cite{W15} are the starting point here. The comparison argument uses
Strassen's coupling theorem~\cite{Str65}, the convex-order criteria in Shaked and Shanthikumar~\cite{SS07} and the
rearrangement-invariant norms of Bennett and Sharpley~\cite{BS88}. Yin and Zhu~\cite{YZ} already gave symmetric laws
that are not $3$-completely mixable, and the polynomial example below is a positive-density refinement of theirs. We
identify Wang's norm condition for bounded laws with a convex order, give non-symmetric counterexamples for $n=4$, and
determine the Bernoulli norm range.

\begin{definition}[{\cite[Definition~5]{W15}}]\label{def:norm}
A \emph{law-determined norm} is a map $N\colon L^0\to[0,\infty]$ such that for all $X,Y\in L^0$ and $a\in\R$:
(i) $N(aX)=|a|N(X)$; (ii) $N(X+Y)\le N(X)+N(Y)$; (iii) $N(X)=0$ implies $X=0$; (iv) $N(X)=N(Y)$ if $X\eqd Y$;
(v) $N(X)\le N(Y)$ if $0\le X\le Y$.
\end{definition}

The value $+\infty$ is allowed, and neither completeness nor continuity is assumed. Write $x_+=\max(x,0)$ and
$x_-=\max(-x,0)$.

\begin{theorem}[{Wang~\cite[Theorem~5]{W15}}]\label{thm:wang}
Let $F$ be $n$-CM with mean $\mu$ and $X\sim F$. Then for every law-determined norm $N$ and every $t\in\R$, with
$s=(n\mu-t)/(n-1)$,
\begin{equation}\label{eq:Nn}
N\bigl((X-t)_+\bigr)\le(n-1)N\bigl((X-s)_-\bigr),\qquad N\bigl((X-t)_-\bigr)\le(n-1)N\bigl((X-s)_+\bigr).
\end{equation}
\end{theorem}

We say that $F$ satisfies $\cN_n$ if~\eqref{eq:Nn} holds for every law-determined norm and every $t$. Open
Problem~4 of~\cite{W15} asks under which additional conditions, for instance smoothness of the density, $\cN_n$ is
sufficient for $n$-complete mixability, and asks for a characterization of law-determined norms. Our first result
removes the norms from the question.

\begin{theorem}\label{thm:A}
Let $X$ have finite mean $\mu$, let $n\ge2$ and put $Y=n\mu-(n-1)X$.
\begin{enumerate}
\item If $X$ is bounded, $X$ satisfies $\cN_n$ if and only if
\begin{equation}\label{eq:cx}
X\cx Y,\qquad\text{equivalently}\qquad X-\mu\cx-(n-1)(X-\mu).
\end{equation}
\item If $X$ is integrable, $\cN_n$ implies~\eqref{eq:cx}, and~\eqref{eq:cx} implies~\eqref{eq:Nn} for every
law-determined norm with the Fatou property~\eqref{eq:fatou}.
\item If $X$ is integrable, $\cN_2$ holds if and only if $F$ is symmetric about $\mu$, which holds if and only if $F$
is $2$-CM.
\end{enumerate}
\end{theorem}

Here $\cx$ is the convex order: $U\cx V$ if $\E\phi(U)\le\E\phi(V)$ for every convex $\phi$ for which both
expectations exist. Proposition~\ref{prop:nonfatou} shows that some restriction of this kind is needed in part (2)
for the norm comparison underlying the proof: for unbounded variables, weak majorization does not imply comparison
in every law-determined norm. The equivalence of $2$-complete mixability with symmetry is due to Wang and
Wang~\cite[Proposition~2.3]{WW11}; part~(3) adds that $\cN_2$ already forces symmetry.

With Theorem~\ref{thm:A}, a counterexample to the converse of Theorem~\ref{thm:wang} is a law that satisfies the
convex-order condition~\eqref{eq:cx} and is not $n$-CM. Every bounded law symmetric about its mean satisfies
\eqref{eq:cx} (Lemma~\ref{lem:sym}), so symmetric laws that are not $3$-CM, such as those in Examples~2.1 and~2.3 of
Yin and Zhu~\cite{YZ}, already give counterexamples for $n=3$. The following results show that neither a smooth
density, nor a density bounded away from zero, nor the absence of symmetry rescues the converse.

\begin{theorem}\label{thm:B}
$\cN_n$ does not imply $n$-complete mixability for $n=3,4$, even for laws on $[0,1]$ with a density that is bounded
above and bounded away from zero. For $n=3$ the analytic density
\begin{equation}\label{eq:poly}
f(x)=\frac1{10}+\frac{27}{10}(2x-1)^2,\qquad0\le x\le1,
\end{equation}
is such an example. For $n=3$ and $n=4$ there are also $C^\infty$ densities on $\R$, supported on $[0,1]$ and positive on
$(0,1)$, that satisfy $\cN_n$ and are not $n$-CM; for $n=4$ the examples are not symmetric.
\end{theorem}

\begin{theorem}\label{thm:C}
Let $0<p<1$ and $n\ge2$. The Bernoulli$(p)$ law satisfies $\cN_n$ if and only if $1/n\le p\le1-1/n$. It is $n$-CM
if and only if $np$ is an integer. In particular Bernoulli$(1/2)$ satisfies $\cN_3$ and is not $3$-CM.
\end{theorem}

The second statement of Theorem~\ref{thm:C} is known: Wang and Wang~\cite[Proposition~2.3]{WW11} construct the
complete mixes, and the set of attainable laws of a sum of Bernoulli variables is determined by Mao and
Wang~\cite[Example~3.1]{MW15}; a related divisibility criterion for binary matrices is due to Haus~\cite{Haus15}, as recalled
in~\cite[Theorem~1]{BKP}. The new part is the norm criterion. For laws with finite support, complete
mixability is decided by the count-vector condition of Puccetti, Wang and Wang~\cite[Theorem~3.2(b)]{PWW12}, which
we recall as Proposition~\ref{prop:pww} and compare with $\cN_n$ on an example.

On the positive side, $\cN_n$ is sufficient within some shape classes. For a monotone density on a bounded interval
$[a,b]$, $\cN_n$ implies the mean condition $a+(b-a)/n\le\mu\le b-(b-a)/n$, which is sufficient by Wang and
Wang~\cite[Theorem~2.4]{WW11}; a symmetric unimodal density is $n$-CM for every $n$ by Rüschendorf and
Uckelmann~\cite{RU02}; a concave density on a bounded interval is $n$-CM for every $n\ge3$ by~\cite[Theorem~4.2]{PWW12}; and for $n\ge3$ a density
bounded below by $3/(n(b-a))$ on its support $[a,b]$ is $n$-CM by~\cite[Theorem~3.4]{PWW13}. The examples of Theorem~\ref{thm:B} show
that an arbitrary positive lower bound in place of $3/(n(b-a))$ does not suffice. For other necessary conditions and
numerical tests see Wang and Wang~\cite{WW16}, Puccetti and Wang~\cite{PW15} and Komisarski and Labuschagne~\cite{KL}.

For a fixed law and large $n$ the picture is different. It is proved in~\cite[Theorem~1.2]{ShLarge} that a law supported in an
interval of length $R$ that dominates $q$ times the uniform law on an interval of length $2H$ is $n$-CM for every
\begin{equation}\label{eq:large}
n\ge\frac{6R/H+1}{q}.
\end{equation}
The density~\eqref{eq:poly} and the density of the law $F_3$ of Section~\ref{sec:dens} are at least $1/10$ on
$[0,1]$, so~\eqref{eq:large} with $R=1$, $H=1/2$ and $q=1/10$ makes them $n$-CM for every $n\ge130$. The proof
of~\eqref{eq:large} absorbs the error of the other copies with the uniform pieces of many copies, and for $n$ below
the threshold it gives nothing. The examples here show that at $n=3$ a lower bound of $1/10$ on the density, and at
$n=4$ a lower bound of $43/1000$, cannot take its place, even together with all of Wang's norm inequalities.

Section~\ref{sec:norms} proves the comparison lemma for norms and the example of Proposition~\ref{prop:nonfatou};
Section~\ref{sec:reduction} proves Theorem~\ref{thm:A}; Section~\ref{sec:atoms} treats atoms and proves
Theorem~\ref{thm:C}; Section~\ref{sec:dens} proves Theorem~\ref{thm:B}; Section~\ref{sec:remarks} makes some remarks
on Definition~\ref{def:norm} and states what remains open. Symbols are listed in Appendix~\ref{app:notation}.

\section{Weak majorization and law-determined norms}\label{sec:norms}

For $Z\ge0$ let $Z^*$ be its decreasing rearrangement on $(0,1)$ and $K_Z(a)=\int_0^aZ^*(u)\,du$ for $0\le a\le1$.
For nonnegative integrable $U,V$, the relation $K_U\le K_V$ is weak submajorization, and it is equivalent to the
increasing convex order $U\icx V$~\cite[Section~4.A]{SS07}. A law-determined norm $N$ has the \emph{Fatou property} if
\begin{equation}\label{eq:fatou}
0\le Z_j\uparrow Z\quad\Longrightarrow\quad N(Z_j)\uparrow N(Z),
\end{equation}
infinite limits included. Definition~\ref{def:norm} does not require~\eqref{eq:fatou}.

\begin{lemma}\label{lem:comp}
Let $U,V\ge0$ be integrable.
\begin{enumerate}
\item The following are equivalent: (a) $K_U(a)\le K_V(a)$ for $0\le a\le1$; (b) $\E(U-z)_+\le\E(V-z)_+$ for $z\ge0$;
(c) $N(U)\le N(V)$ for every law-determined norm with the Fatou property.
\item If $U$ and $V$ are bounded, (a) is equivalent to $N(U)\le N(V)$ for every law-determined norm.
\end{enumerate}
\end{lemma}

\begin{proof}
The identities
\[
K_Z(a)=\inf_{z\ge0}\{az+\E(Z-z)_+\},\qquad\E(Z-z)_+=\sup_{0\le a\le1}\{K_Z(a)-az\}
\]
follow by splitting $\int_0^aZ^*$ at a point where $Z^*$ crosses $z$, and they give (a)$\Leftrightarrow$(b). For
$0<a\le1$ the functional
\[
T_a(Z)=\frac1a\int_0^a|Z|^*(u)\,du=\frac1a\sup\{\E[|Z|H]:0\le H\le1,\ \E H\le a\}
\]
is a law-determined norm with the Fatou property: the dual formula gives subadditivity and monotone convergence gives
\eqref{eq:fatou}. Hence comparison in all norms of either class implies (a).

For the converse let $N$ be a law-determined norm and $\nu(Z)=N(|Z|)$. Then $\nu$ agrees with $N$ on nonnegative
variables, and it is again a law-determined norm, since $|Z+W|\le|Z|+|W|$ together with (v) and (ii) gives
subadditivity. If $\nu(V)=\infty$ there is nothing to prove, and if $V=0$ then (a) forces $U=0$. Otherwise pick
$\varepsilon>0$ with $p=\Prob(V\ge\varepsilon)>0$; by (v) the indicator of an event of probability $p$ has finite
norm. Finitely many translates of an arc of length $p$ cover the unit circle, so finitely many indicators of events of
probability $p$ sum to at least $1$, and $C=\nu(1)<\infty$. Consequently, for bounded $Z,W$,
\begin{equation}\label{eq:lip}
\nu(Z)\le C\|Z\|_\infty,\qquad|\nu(Z)-\nu(W)|\le C\|Z-W\|_\infty.
\end{equation}

We show that $\nu(\E[V\mid\mathcal P])\le\nu(V)$ for bounded $V\ge0$ and every finite partition $\mathcal P$. On a cell
of positive probability let $T$ be uniform on $(0,1)$ and $q$ the conditional quantile function of $V$, and put
$V_j=q((T+j/m)\bmod1)$ on that cell, $0\le j<m$. Each $V_j$ has the law of $V$, and if $0\le V\le M$ the average
$m^{-1}\sum_jV_j$, an equally spaced Riemann sum of the monotone function $q$, is within $M/m$ of
$\E[V\mid\mathcal P]$ uniformly. Subadditivity, (iv) and~\eqref{eq:lip} give the claim as $m\to\infty$.

By Strassen's theorem~\cite[Theorem~9]{Str65} (in the form for the increasing convex order,
see~\cite[Theorem~4.A.5]{SS07}), (a) gives a coupling with $U\le\E[V\mid U]$ almost surely. Suppose first that $U$ and
$V$ are bounded. Partition the range of $U$ into intervals of length at most $\delta$ and let $U_\delta$ be the left
endpoint of the interval containing $U$. Then $U_\delta\le\E[V\mid\mathcal P_\delta]$, where $\mathcal P_\delta$ is the
partition generated by $U_\delta$, so by (v) and the previous paragraph $\nu(U_\delta)\le\nu(V)$, and
by~\eqref{eq:lip} $\nu(U)\le\nu(V)+C\delta$. Letting $\delta\to0$ proves (2).

If $U,V$ are integrable and $N$ has the Fatou property, put $W_j=\E[V\wedge j\mid U]$. Conditional Jensen gives
$W_j\cx V\wedge j$, so the bounded case and (v) give $\nu(W_j)\le\nu(V\wedge j)\le\nu(V)$. By conditional monotone
convergence $W_j\uparrow\E[V\mid U]\ge U$, and~\eqref{eq:fatou} and (v) give $\nu(U)\le\nu(V)$.
\end{proof}

The Fatou property cannot be dropped from part (1) of Lemma~\ref{lem:comp}.

\begin{proposition}\label{prop:nonfatou}
There are nonnegative integrable $U,V$ with $U\cx V$ and a law-determined norm $N$ with $N(V)<\infty=N(U)$.
\end{proposition}

\begin{proof}
On $(0,1)$ let $v(t)=(t\log^2(1/t))^{-1}$ for $0<t<e^{-3}$ and $v(t)=0$ for $e^{-3}\le t<1$; $v$ is decreasing with
integral $1/3$. Let $T$ be uniform on $(0,1)$ and $V=v(T)$. For $k=1,2,\dots$ put $L_k=3^k$, $a_k=e^{-2L_k}$,
$b_k=e^{-L_k}$ and $I_k=(a_k,b_k)$; these intervals are disjoint and lie in $(0,e^{-3})$, and
$\int_{I_k}v=1/(2L_k)$. Let $U=u(T)$, where $u=v$ off $\bigcup_kI_k$ and $u$ equals the average
$h_k=(2L_k(b_k-a_k))^{-1}$ of $v$ on $I_k$. Then $U$ is a conditional expectation of $V$, so $U\cx V$.

Let $D$ be the set of $Z\in L^0$ such that $|Z|^*(t)\le Cv(t/m)$ for almost every $t\in(0,1)$, for some $C<\infty$
and some integer $m\ge1$. The inequality $|Z+W|^*(t)\le|Z|^*(t/2)+|W|^*(t/2)$ shows that $D$ is a linear space; it
is solid and law invariant, it is contained in $L^1$, and it contains $V$ and all bounded variables (take $m>e^3$).
Define $N(Z)=\E|Z|$ for $Z\in D$ and $N(Z)=\infty$ otherwise. The properties of $D$ give (i)--(v).

Put $t_k=(b_k-a_k)/2$. Since $u=h_k$ on a set of measure $2t_k$, $U^*(t)\ge h_k$ for $t\le t_k$, and for every fixed
$m$ and every $t\in[t_k/2,t_k]$,
\[
\frac{U^*(t)}{v(t/m)}\ge\frac{h_k}{v(t_k/(2m))}=\frac{t_k\log^2(2m/t_k)}{2m\cdot2L_k(b_k-a_k)}
=\frac{\log^2(2m/t_k)}{8mL_k}\to\infty
\]
as $k\to\infty$, because $\log(1/t_k)\ge L_k$. Hence $U\notin D$ and $N(U)=\infty$, while $N(V)=\E V=1/3$. (Also
$N(U\wedge j)=\E[U\wedge j]\uparrow\E U$, so $N$ fails~\eqref{eq:fatou}.)
\end{proof}

Proposition~\ref{prop:nonfatou} concerns the comparison lemma, not Theorem~\ref{thm:wang}: we do not know an
integrable law $F$ that satisfies~\eqref{eq:cx} and violates~\eqref{eq:Nn} for some law-determined norm without the
Fatou property.

\section{Reduction to a convex order; the case \texorpdfstring{$n=2$}{n=2}}\label{sec:reduction}

\begin{proposition}\label{prop:red}
Parts (1) and (2) of Theorem~\ref{thm:A} hold. Moreover, for bounded $X$ the first inequality in~\eqref{eq:Nn} for the
norm $\E|\cdot|$ and all $t$ already implies $\cN_n$.
\end{proposition}

\begin{proof}
Put $c=n-1$ and $s=(n\mu-t)/c$. Then pointwise
\begin{equation}\label{eq:hinge}
c(X-s)_-=(Y-t)_+,\qquad c(X-s)_+=(Y-t)_-.
\end{equation}
So the first inequality of~\eqref{eq:Nn} for $N=\E|\cdot|$ reads $\E(X-t)_+\le\E(Y-t)_+$ for all $t$, which, since
$\E X=\E Y=\mu$, is $X\cx Y$~\cite[Theorem~3.A.1]{SS07}. Conversely, if $X\cx Y$, then for fixed $t$ both $x\mapsto(x-t)_+$
and $x\mapsto(t-x)_+$ are convex, and composing them with increasing convex functions gives
\[
(X-t)_+\icx(Y-t)_+,\qquad(X-t)_-\icx(Y-t)_-.
\]
Lemma~\ref{lem:comp} and~\eqref{eq:hinge} give~\eqref{eq:Nn}, for every law-determined norm if $X$ is bounded and
for every law-determined norm with the Fatou property if $X$ is integrable.
\end{proof}

The condition~\eqref{eq:cx} is a consequence of $n$-complete mixability that can be seen directly: if $X_1,\dots,X_n$
is a complete mix and $Z_i=X_i-\mu$, then $Z_1=-\sum_{i\ge2}Z_i$, and for convex $\phi$,
\[
\E\phi(Z_1)\le\frac1{n-1}\sum_{i=2}^n\E\phi\bigl(-(n-1)Z_i\bigr)=\E\phi\bigl(-(n-1)Z_1\bigr).
\]

\begin{proposition}\label{prop:n2}
For every integrable $F$, $\cN_2$, symmetry of $F$ about its mean, and $2$-complete mixability are equivalent.
\end{proposition}

\begin{proof}
By Proposition~\ref{prop:red}, $\cN_2$ gives $X\cx2\mu-X$. Reflection $x\mapsto2\mu-x$ preserves the convex order, so
also $2\mu-X\cx X$, and the call functions $t\mapsto\E(X-t)_+$ and $t\mapsto\E(2\mu-X-t)_+$ coincide. Their right
derivatives are $-\Prob(X>t)$ and $-\Prob(2\mu-X>t)$, so $X\eqd2\mu-X$. If $F$ is symmetric, $(X,2\mu-X)$ is a
complete mix, and with $s=2\mu-t$ the law of $(X-t)_+$ is that of $(X-s)_-$ and the law of $(X-t)_-$ is that of
$(X-s)_+$; thus every inequality in~\eqref{eq:Nn} is an equality. Finally a complete mix $(X_1,X_2)$ has
$X_2=2\mu-X_1$, so $F$ is symmetric~\cite[Proposition~2.3]{WW11}.
\end{proof}

\section{Atoms}\label{sec:atoms}

\begin{lemma}\label{lem:sym}
Let $Z$ be integrable and symmetric about $0$, and $c\ge1$. Then $Z\cx cZ\eqd-cZ$. Consequently every bounded law
symmetric about its mean satisfies $\cN_n$ for every $n\ge2$.
\end{lemma}

\begin{proof}
Write $Z=c^{-1}(cZ)+(1-c^{-1})\cdot0$. For convex $\phi$, $\phi(Z)\le c^{-1}\phi(cZ)+(1-c^{-1})\phi(0)$, and
$\phi(0)\le\E\phi(cZ)$ by Jensen. The second statement follows from Theorem~\ref{thm:A}(1) with $c=n-1$.
\end{proof}

\begin{proof}[Proof of Theorem~\ref{thm:C}]
Let $B\sim\mathrm{Bernoulli}(p)$ and $Y=np-(n-1)B$, which takes the values $l=np-(n-1)$ and $h=np$. If $B\cx Y$, the
support of $Y$ spans an interval containing $\{0,1\}$ (otherwise a hinge function beyond an endpoint violates the
order), so $np\ge1$ and $np-(n-1)\le0$. Conversely, if $l\le0$ and $h\ge1$, every convex $\phi$ on $[l,h]$ lies below
its chord through $(l,\phi(l))$ and $(h,\phi(h))$; taking expectations at $B$, and using that $Y$ lives on
$\{l,h\}$ and has the same mean as $B$, gives $\E\phi(B)\le\E\phi(Y)$. Theorem~\ref{thm:A}(1) turns this into the
norm statement.

A sum of Bernoulli variables is an integer, so a complete mix requires $np\in\Z$. If $np=k$, the indicators of a
uniformly random $k$-element subset of $\{1,\dots,n\}$ are Bernoulli$(p)$ and sum to $k$~\cite[Proposition~2.3]{WW11}.
\end{proof}

Not every lattice law with an unreachable target passes $\cN_n$: Bernoulli$(1/10)$ fails it for $n=3$. If $F$ is
supported on $a+h\Z$ with $h>0$, an $n$-complete mix needs $n\mu\in na+h\Z$, and more precisely $n\mu\in A+\dots+A$
($n$ summands) for the support $A$. Even reachability, together with the requirement that every atom occurs in some
admissible $n$-tuple, does not suffice; the masses matter. For finite support the exact condition is the following.

\begin{proposition}[{Puccetti, Wang and Wang~\cite[Theorem~3.2(b)]{PWW12}}]\label{prop:pww}
Let $F\{a_j\}=p_j>0$ for $1\le j\le m$, and let $\mathcal K$ be the set of $k\in\Z_{\ge0}^m$ with $\sum_jk_j=n$ and
$\sum_jk_ja_j=n\mu$. Then $F$ is $n$-CM if and only if $(p_1,\dots,p_m)$ lies in the convex hull of
$\{k/n:k\in\mathcal K\}$.
\end{proposition}

In~\cite{PWW12} the statement is that a discrete law is $n$-CM if and only if it is a mixture of empirical laws of
$n$ points with mean $\mu$; on a finite support the empirical law of such an $n$-tuple is the vector $k/n$. The proof is
short: symmetrizing a complete mix by an independent uniform permutation produces a random count vector in
$\mathcal K$ with mean $(np_j)_j$, and conversely a mixture of count vectors followed by a uniform permutation of the
resulting multiset is a complete mix.

\begin{example}\label{ex:three}
Give the points $0,1,2$ the masses $49/100$, $1/50$, $49/100$, and let $n=3$. The law is symmetric about $1$, so it
satisfies $\cN_3$ by Lemma~\ref{lem:sym}. The target $3$ is reachable and every atom occurs in an admissible triple,
but the admissible unordered triples are $(0,1,2)$ and $(1,1,1)$, both containing the point $1$. By
Proposition~\ref{prop:pww} a complete mix would need $\Prob(X=1)\ge1/3$, whereas $\Prob(X=1)=1/50$.
\end{example}

\section{Densities}\label{sec:dens}

The densities below are excluded from complete mixability by a mass argument; its symmetric version is Example~2.1
of Yin and Zhu~\cite{YZ}.

\begin{lemma}\label{lem:mass}
Let $X\sim F$ have mean $\mu$, let $A\subset\R$ be measurable, let $nA=\{x_1+\dots+x_n:x_i\in A\}$, and suppose that
the distance from $n\mu$ to $nA$ is at least $d>0$. Let $q=\Prob(X\notin A)$. Then every coupling $X_1,\dots,X_n$ of
$n$ copies of $F$ satisfies
\begin{equation}\label{eq:mass}
\E\Big(\sum_{i=1}^nX_i-n\mu\Big)^2\ge d^2(1-nq)_+.
\end{equation}
In particular $F$ is not $n$-CM if $q<1/n$.
\end{lemma}

\begin{proof}
With probability at least $1-nq$ every $X_i$ lies in $A$, and then the sum is at distance at least $d$ from $n\mu$.
\end{proof}

\subsection{Two uniform components}
For $0<\varepsilon<1/2$ and $0<p<1$ let
\[
G_{p,\varepsilon}=(1-p)\,\mathrm{U}[0,\varepsilon]+p\,\mathrm{U}[1-\varepsilon,1],
\]
the law of $(1-\varepsilon)B+\varepsilon U$ with independent $B\sim\mathrm{Bernoulli}(p)$ and $U\sim\mathrm U[0,1]$. Its
mean is $(1-\varepsilon)p+\varepsilon/2$.

\begin{lemma}\label{lem:G}
If $1/n\le p\le1-1/n$, then $G_{p,\varepsilon}$ satisfies $\cN_n$.
\end{lemma}

\begin{proof}
By the proof of Theorem~\ref{thm:C}, $B-p\cx-(n-1)(B-p)$, and by Lemma~\ref{lem:sym} the same holds for $U-1/2$. The
convex order is preserved under sums of independent variables (condition on one summand, then on the
other~\cite[Theorem~3.A.12]{SS07}), so it holds for $(1-\varepsilon)(B-p)+\varepsilon(U-1/2)$; apply
Theorem~\ref{thm:A}(1).
\end{proof}

Let $G_3=G_{1/2,1/10}$ and $G_4=G_{2/5,1/20}$, with densities
\[
g_3=5\cdot\one_{[0,1/10]\cup[9/10,1]},\qquad g_4=12\cdot\one_{[0,1/20]}+8\cdot\one_{[19/20,1]}.
\]
Both satisfy $\cN_3$ and $\cN_4$ respectively. For $A_\varepsilon=[0,\varepsilon]\cup[1-\varepsilon,1]$,
\begin{equation}\label{eq:nA}
nA_\varepsilon=\bigcup_{k=0}^n\big[k(1-\varepsilon),\,k(1-\varepsilon)+n\varepsilon\big].
\end{equation}
For $G_3$ the target $3/2$ lies between the intervals ending at $6/5$ and starting at $9/5$, at distance $3/10$ from
both. For $G_4$ the mean is $81/200$ and the target $81/50$ lies between the intervals ending at $23/20$ and starting
at $19/10$, at distances $47/100$ and $7/25$. Since $\Prob(X\in A_\varepsilon)=1$, Lemma~\ref{lem:mass} with $q=0$
shows that $G_3$ is not $3$-CM and $G_4$ is not $4$-CM.

\subsection{A positive density}
Let
\begin{equation}\label{eq:F3F4}
F_3=\tfrac9{10}G_3+\tfrac1{10}\mathrm U[0,1],\qquad F_4=\tfrac9{10}G_4+\tfrac1{10}H_4,
\end{equation}
where $H_4$ has density $h_4(x)=\frac{157}{100}-\frac{57}{50}x$ on $[0,1]$. The density $h_4$ integrates to $1$, is at
least $43/100$, and $H_4$ has mean $81/200$.

\begin{proposition}\label{prop:F}
$F_3$ satisfies $\cN_3$ and is not $3$-CM; $F_4$ satisfies $\cN_4$ and is not $4$-CM. Their densities are bounded below
on $[0,1]$ by $1/10$ and $43/1000$.
\end{proposition}

\begin{proof}
$F_3$ is symmetric about $1/2$, so Lemma~\ref{lem:sym} applies. For $F_4$ we check $H_4\cx Y_H$, where
$Y_H=4\mu-3H_4$ and $\mu=81/200$. For a law on $[0,1]$ with mean $\mu$ it suffices to compare call functions on
$[0,1]$: for $t\le0$ the call function of $H_4$ is $\mu-t$, which by Jensen is at most that of $Y_H$, and for $t\ge1$ it
is $0$. On $[0,1]$ the call function of $H_4$ is $C(t)=(1-t)^2(81-38t)/200$, and since $s=(4\mu-t)/3\in[0,1]$ for
$t\in[0,1]$, the call function of $Y_H$ at $t$ is $3\,\E(s-H_4)_+=3\{C(s)+s-\mu\}$. Their difference is
\[
D(t)=3\{C(s)+s-\mu\}-C(t)=\frac{11875000t^3-35208750t^2+17910675t+10797948}{56250000},
\]
and its Bernstein coefficients of degree three on $[0,1]$ are
\[
\frac{299943}{1562500},\quad\frac{5589391}{18750000},\quad\frac{2750537}{14062500},\quad\frac{5374873}{56250000}.
\]
They are positive, so $D\ge0$ on $[0,1]$ and $H_4\cx Y_H$. The convex-order inequality~\eqref{eq:cx} is linear in the
law when the mean is fixed, so it passes to the mixture $F_4$ of $G_4$ (Lemma~\ref{lem:G}) and $H_4$, both of mean
$81/200$; Theorem~\ref{thm:A}(1) gives $\cN_4$.

For $F_3$ take $A=[0,1/10]\cup[9/10,1]$: $q=\frac1{10}\cdot\frac45=\frac2{25}$ and $d=3/10$, so~\eqref{eq:mass} gives
$\E(S-3/2)^2\ge\frac9{100}(1-\frac6{25})=\frac{171}{2500}$ for every coupling, where $S$ is the sum. For $F_4$ take
$A=[0,1/20]\cup[19/20,1]$: integrating $h_4$ over $[1/20,19/20]$ gives $9/10$, so $q=9/100$, and~\eqref{eq:mass} gives
$\E(S-81/50)^2\ge\frac{49}{625}(1-\frac{36}{100})=\frac{784}{15625}$. The density bounds come from the uniform and
$H_4$ components.
\end{proof}

\subsection{An analytic density and smooth densities}

\begin{proposition}\label{prop:poly}
The density~\eqref{eq:poly} satisfies $\cN_3$ and is not $3$-CM.
\end{proposition}

\begin{proof}
It integrates to $1$ and is symmetric about $1/2$, so Lemma~\ref{lem:sym} applies. If $x_1,x_2,x_3\in[0,1]$ sum to
$3/2$, one of them lies in $[1/4,3/4]$: otherwise the numbers $2x_i-1$ all have absolute value greater than $1/2$, two
of them have the same sign, and their sum has absolute value greater than $1$, which the third cannot cancel. So a
complete mix would have at least one coordinate in $[1/4,3/4]$ almost surely, and $3\Prob(X\in[1/4,3/4])\ge1$. But
\[
\Prob(X\in[1/4,3/4])=\frac1{10}\cdot\frac12+\frac9{10}\cdot\Big(\frac12\Big)^3=\frac{13}{80}<\frac13.\qedhere
\]
\end{proof}

The component $3(2x-1)^2$ of~\eqref{eq:poly} is an affine image of the case $r=1$ of Example~2.3 of Yin and
Zhu~\cite{YZ}; the uniform component makes the density positive on the closed interval without removing the
obstruction. Lemma~\ref{lem:mass} with $A=[0,1/5]\cup[4/5,1]$, $d=1/10$ and $q=159/625$ gives the quantitative bound
$\E(S-3/2)^2\ge37/15625$.

\begin{corollary}\label{cor:smooth}
For $n=3,4$ there is a $C^\infty$ density on $\R$, supported on $[0,1]$ and positive on $(0,1)$, whose law satisfies
$\cN_n$ and is not $n$-CM. For $n=4$ it is not symmetric.
\end{corollary}

\begin{proof}
Let $\eta=1/100$ and let $W$ be independent of $X\sim F_n$, with density proportional to
$\exp(-1/(1-(w/\eta)^2))$ on $(-\eta,\eta)$ and $0$ elsewhere. The law of $X+W$ has a $C^\infty$ density, supported on
$[-\eta,1+\eta]$ and positive inside. Its centered convex order~\eqref{eq:cx} follows from those of $X$ and $W$
(Lemma~\ref{lem:sym}) by independence, as in Lemma~\ref{lem:G}. With probability at least $9/10$ the variable $X+W$ lies
in the enlarged core $A+[-\eta,\eta]$, $A=A_{1/10}$ or $A_{1/20}$, and the distance from the target to its $n$-fold sum
set is at least $d_n-n\eta$, with $d_3=3/10$ and $d_4=7/25$. Since $1/10<1/n$, Lemma~\ref{lem:mass} excludes a complete
mix. The affine map $x\mapsto(x+\eta)/(1+2\eta)$ moves the support to $[0,1]$ and preserves both properties; for
$n=4$ the mean becomes $83/204\ne1/2$.
\end{proof}

A $C^\infty$ density with compact support tends to $0$ at the ends of its support, so Corollary~\ref{cor:smooth} and
Proposition~\ref{prop:F} give different kinds of regularity; the density~\eqref{eq:poly} has both a positive lower
bound and analyticity on the closed interval. None of this contradicts the sufficient lower bound $3/(n(b-a))$
of~\cite[Theorem~3.4]{PWW13}: for $n=3$ on $[0,1]$ it equals $1$, and the density~\eqref{eq:poly} dips to $1/10$.

\begin{proof}[Proof of Theorem~\ref{thm:B}]
Propositions~\ref{prop:F} and~\ref{prop:poly} and Corollary~\ref{cor:smooth}.
\end{proof}

\section{Remarks on \texorpdfstring{Definition~\ref{def:norm}}{the definition of a norm}, and open questions}\label{sec:remarks}

Every argument in~\eqref{eq:Nn} is nonnegative, so only the restriction of a norm to nonnegative variables enters.
These restrictions are exactly the maps $\nu$ from nonnegative variables to $[0,\infty]$ that are positively
homogeneous, subadditive, definite, law invariant and monotone, since each such $\nu$ extends to a law-determined norm
by $Z\mapsto\nu(|Z|)$ (the proof of Lemma~\ref{lem:comp}). The definition itself is broader than the absolute
rearrangement-invariant norms: for $\alpha>0$,
\[
N_\alpha(Z)=\E|Z|+\alpha|\E Z|\ \ (Z\in L^1),\qquad N_\alpha(Z)=\infty\ \ (Z\notin L^1),
\]
satisfies (i)--(v), and a symmetric $\pm1$ variable $Z$ has $N_\alpha(Z)=1<1+\alpha=N_\alpha(|Z|)$. Among absolute norms
with the Fatou property that are finite on bounded variables, one recovers the rearrangement-invariant Banach function
norms on $(0,1)$, which have the Lorentz--Luxemburg representation as suprema of weighted integrals
$\int_0^1|Z|^*(u)w(u)\,du$ over decreasing weights~\cite[Chapter~2]{BS88}; Proposition~\ref{prop:nonfatou} shows that
Definition~\ref{def:norm} also admits norms outside this class.

For bounded laws, Theorem~\ref{thm:A} says that no choice of law-determined norms gives more than~\eqref{eq:cx}, so
Open Problem~4 becomes a question about the convex order: for which laws does $X-\mu\cx-(n-1)(X-\mu)$ imply
$n$-complete mixability? The following questions remain open.
\begin{enumerate}
\item Does~\eqref{eq:cx} imply $n$-complete mixability for unimodal densities that are not symmetric?
\item For fixed $n$ and mean, what is the smallest lower bound on a density on $[0,1]$ that, together with
\eqref{eq:cx}, implies $n$-complete mixability? At mean $81/200$ for $n=4$ the examples above show that it is at
least $43/1000$, and at mean $1/2$ for $n=3$ at least $1/10$, among laws satisfying~\eqref{eq:cx}, and~\cite[Theorem~3.4]{PWW13} shows that $3/n$ suffices
without~\eqref{eq:cx}.
\item For unbounded integrable laws, does~\eqref{eq:cx} imply~\eqref{eq:Nn} for every law-determined norm, including
those without the Fatou property?
\end{enumerate}
The counterexamples are for fixed $n$ and say nothing about large $n$ for a fixed law; compare~\cite{ShLarge}.

\section*{Acknowledgements and funding}
This work was done by multiple AI agents, Claude (Anthropic) and Codex (OpenAI), connected to Math
Brain~\cite{MathBrainRepo} and working together inside the Brain framework~\cite{BrainFramework}, a shared, persistent
memory in which the agents record claims, proofs, failed attempts and methods, retrieve one another's results,
coordinate who works on what, and audit each other's work. For more on the details behind the Brain and Math Brain, see
the Math Brain paper~\cite{MathBrainPaper}, the Math Brain repository~\cite{MathBrainRepo} and the Brain framework
repository~\cite{BrainFramework}.

The research presented in this paper was supported by the European Research Council (ERC) under the European Union's
Horizon Europe research and innovation programme (grant agreement No.~101041711), by the Simons Foundation as part of
the Collaboration on the Mathematical and Scientific Foundations of Deep Learning, by Heights Labs, by Convex Nexus
Capital, by the Israel Science Foundation (grant number 2258/19), and by the Israel Science Foundation (ISF Grant
4101/25). Computational resources were provided in part by Google Cloud through its education credits program (credit
530550103).

\section*{Declaration of competing interest}
The author holds equity in Heights Labs and in Convex Nexus Capital and has a management role at Convex Nexus Capital.
These entities are included in the funding declaration above.

\section*{Declaration of generative AI and AI-assisted technologies}
The paper comes out of Math Brain, an ongoing project on the use of AI systems for research mathematics, and its
results were developed with the Brain~\cite{BrainFramework} and Math Brain~\cite{MathBrainRepo}. A Codex agent
(OpenAI) developed the norm reduction and the Bernoulli norm criterion, constructed the examples and performed the
exact checks; a second Codex agent checked the proofs and recomputed every certificate with its own program; a Claude
agent (Anthropic) prepared the manuscript. The results used from the literature are those of their authors: Ruodu
Wang's norm inequalities and question~\cite{W15}; Bin Wang and Ruodu Wang's symmetry and Bernoulli results~\cite{WW11};
Tiantian Mao and Ruodu Wang's description of Bernoulli aggregation~\cite{MW15}; the discrete criterion and density
conditions of Giovanni Puccetti, Bin Wang and Ruodu Wang~\cite{PWW12,PWW13}; the symmetric unimodal construction of
Ludger R\"uschendorf and Ludger Uckelmann~\cite{RU02}; the symmetric non-mixable laws of Chuancun Yin and Dan
Zhu~\cite{YZ}; Volker Strassen's coupling theorem~\cite{Str65}; and the stochastic-order and rearrangement-invariant
norm theory in Moshe Shaked and J.~George Shanthikumar~\cite{SS07} and Colin Bennett and Robert Sharpley~\cite{BS88}.
The binary-matrix criterion mentioned in the introduction is due to Utz-Uwe Haus~\cite{Haus15}. New here are the norm reduction, its consequences for the norm condition, and the examples. These checks are not peer
review. The author directed the work and is responsible for the content of the paper and
for decisions concerning its dissemination.

\section*{Code and data availability}
The programs and their outputs are in the repository
\url{https://github.com/yspennstate/complete-mixability-norm-converse/tree/v1.1}. The exact checks use SymPy~\cite{SymPy}; the linear programs use the HiGHS solver~\cite{HiGHS} through
SciPy~\cite{SciPy}.

The paper is part of the Math Brain project, whose public repository is~\cite{MathBrainRepo}.

\begin{thebibliography}{99}
\bibitem{BKP} F.~Bellini, O.~E.~Karaşan and M.~Ç.~Pınar, Joint mixability of some integer matrices, \emph{Discrete
Optim.} \textbf{20} (2016), 90--104.
\url{https://doi.org/10.1016/j.disopt.2016.03.003}
\bibitem{BS88} C.~Bennett and R.~Sharpley, \emph{Interpolation of Operators}, Pure and Applied Mathematics
\textbf{129}, Academic Press, Boston, 1988.
\bibitem{HiGHS} Q.~Huangfu and J.~A.~J.~Hall, Parallelizing the dual revised simplex method, \emph{Math. Program.
Comput.} \textbf{10} (2018), 119--142.
\url{https://doi.org/10.1007/s12532-017-0130-5}
\bibitem{Haus15} U.-U.~Haus, Bounding stochastic dependence, joint mixability of matrices, and multidimensional
bottleneck assignment problems, \emph{Oper. Res. Lett.} \textbf{43} (2015), 74--79.
\bibitem{KL} A.~Komisarski and J.~Labuschagne, The distributions of the mean of random vectors with fixed marginal
distribution, \emph{J. Theoret. Probab.} \textbf{37} (2024), 2121--2129.
\url{https://doi.org/10.1007/s10959-023-01277-2}
\bibitem{MW15} T.~Mao and R.~Wang, On aggregation sets and lower-convex sets, \emph{J. Multivariate Anal.}
\textbf{138} (2015), 170--181.
\url{https://doi.org/10.1016/j.jmva.2014.12.001}
\bibitem{SymPy} A.~Meurer et al., SymPy: symbolic computing in Python, \emph{PeerJ Comput. Sci.} \textbf{3} (2017),
e103.
\url{https://doi.org/10.7717/peerj-cs.103}
\bibitem{PW15} G.~Puccetti and R.~Wang, Detecting complete and joint mixability, \emph{J. Comput. Appl. Math.}
\textbf{280} (2015), 174--187.
\url{https://doi.org/10.1016/j.cam.2014.11.050}
\bibitem{PWW12} G.~Puccetti, B.~Wang and R.~Wang, Advances in complete mixability, \emph{J. Appl. Probab.} \textbf{49}
(2012), 430--440.
\url{https://doi.org/10.1239/jap/1339878796}
\bibitem{PWW13} G.~Puccetti, B.~Wang and R.~Wang, Complete mixability and asymptotic equivalence of worst-possible
VaR and ES estimates, \emph{Insurance Math. Econom.} \textbf{53} (2013), 821--828.
\url{https://doi.org/10.1016/j.insmatheco.2013.09.017}
\bibitem{RU02} L.~Rüschendorf and L.~Uckelmann, Variance minimization and random variables with constant sum, in
\emph{Distributions with Given Marginals and Statistical Modelling}, Kluwer, Dordrecht, 2002, 211--222.
\url{https://doi.org/10.1007/978-94-017-0061-0_22}
\bibitem{SS07} M.~Shaked and J.~G.~Shanthikumar, \emph{Stochastic Orders}, Springer Series in Statistics, Springer,
New York, 2007. \url{https://doi.org/10.1007/978-0-387-34675-5}
\bibitem{ShLarge} Y.~Shmalo, \emph{Absolutely continuous laws on a bounded interval are completely mixable for all large
$n$}, preprint, 2026,
\url{https://github.com/yspennstate/complete-mixability-large-n/tree/4bd67dbdd5b621c08a813b909d10219d4dd2d18f}.
\bibitem{Str65} V.~Strassen, The existence of probability measures with given marginals, \emph{Ann. Math. Statist.}
\textbf{36} (1965), 423--439.
\url{https://doi.org/10.1214/aoms/1177700153}
\bibitem{SciPy} P.~Virtanen et al., SciPy 1.0: fundamental algorithms for scientific computing in Python, \emph{Nat.
Methods} \textbf{17} (2020), 261--272.
\url{https://doi.org/10.1038/s41592-019-0686-2}
\bibitem{WW11} B.~Wang and R.~Wang, The complete mixability and convex minimization problems with monotone marginal
densities, \emph{J. Multivariate Anal.} \textbf{102} (2011), 1344--1360.
\url{https://doi.org/10.1016/j.jmva.2011.05.002}
\bibitem{WW16} B.~Wang and R.~Wang, Joint mixability, \emph{Math. Oper. Res.} \textbf{41} (2016), 808--826.
\url{https://doi.org/10.1287/moor.2015.0755}
\bibitem{W15} R.~Wang, Current open questions in complete mixability, \emph{Probab. Surv.} \textbf{12} (2015),
13--32.
\url{https://doi.org/10.1214/14-PS250}
\bibitem{YZ} C.~Yin and D.~Zhu, Joint mixability of elliptical distributions and related families, preprint, 2017, revised 2023,
\href{https://arxiv.org/abs/1706.05499v4}{arXiv:1706.05499v4}.
\bibitem{MathBrainRepo} Y.~Shmalo, \emph{Math Brain: public snapshot, code and tests}, 2026,
\url{https://github.com/yspennstate/math-brain-public/tree/8c9a6bd478c4ec4082347151b7ae16db50c917dd}.
\bibitem{MathBrainPaper} Y.~Shmalo, \emph{Math Brain: shared, auditable memory for multi-agent mathematical research},
October 2026, \url{https://github.com/yspennstate/math-brain-public/tree/main/papers/math-brain-2026-10-24},
doi:10.5281/zenodo.24150837.
\bibitem{BrainFramework} Y.~Shmalo, \emph{Brain Framework}, public software repository, 2026,
\url{https://github.com/yspennstate/brain-framework-public/tree/5629e644cdd510173ae97e9aa92642dabd22dc57}.
\end{thebibliography}

\appendix
\section{Notation}\label{app:notation}
Symbols are listed in the order in which they first appear. Letters used only inside one proof are not listed.

\medskip
\begin{center}
\small
\begin{tabular}{@{}lp{0.6\textwidth}l@{}}
\toprule
Symbol & Meaning & First appears\\
\midrule
$F$, $\mu$, $n$ & a law on $\R$, its mean, the number of copies & Section~\ref{sec:intro}\\
$n$-CM & $n$-completely mixable & Section~\ref{sec:intro}\\
$L^0$ & all random variables on the atomless probability space & Section~\ref{sec:intro}\\
$N$ & a law-determined norm (Definition~\ref{def:norm}) & Section~\ref{sec:intro}\\
$x_+$, $x_-$ & $\max(x,0)$ and $\max(-x,0)$ & Section~\ref{sec:intro}\\
$t$, $s$ & a threshold and $s=(n\mu-t)/(n-1)$ & Theorem~\ref{thm:wang}\\
$\cN_n$ & \eqref{eq:Nn} for every law-determined norm and every $t$ & Section~\ref{sec:intro}\\
$Y$ & $n\mu-(n-1)X$ & Theorem~\ref{thm:A}\\
$\cx$, $\icx$ & convex order, increasing convex order & Theorem~\ref{thm:A}, Section~\ref{sec:norms}\\
$Z^*$, $K_Z$ & decreasing rearrangement of $Z\ge0$ and its integral & Section~\ref{sec:norms}\\
$T_a$ & the norm $a^{-1}\int_0^a|Z|^*$ & Lemma~\ref{lem:comp}\\
$\nu$ & the absolute extension $Z\mapsto N(|Z|)$ & Lemma~\ref{lem:comp}\\
$\mathcal K$ & admissible count vectors & Proposition~\ref{prop:pww}\\
$nA$ & the $n$-fold sum set of $A$ & Lemma~\ref{lem:mass}\\
$G_{p,\varepsilon}$, $G_3$, $G_4$ & two-component uniform mixtures & Section~\ref{sec:dens}\\
$A_\varepsilon$ & $[0,\varepsilon]\cup[1-\varepsilon,1]$ & \eqref{eq:nA}\\
$F_3$, $F_4$, $H_4$ & laws with positive densities & \eqref{eq:F3F4}\\
$S$ & the sum of a coupling of $n$ copies & Proposition~\ref{prop:F}\\
$N_\alpha$ & a law-determined norm that is not absolute & Section~\ref{sec:remarks}\\
\bottomrule
\end{tabular}
\end{center}

\end{document}
